% ECONOMICS_PROBLEM: {"id": "D63-158", "title": "Ex ante envy-freeness with ex post EF1 and Pareto optimality", "jel_code": "D63", "source_id": "OP-0191"}
% FIRST_POSTED: 2026-10-09
% ATTEMPT_STATUS: unresolved
% ENTRY_KIND: research_attempt
% !TEX program = pdflatex
\documentclass[11pt]{article}
\usepackage[T1]{fontenc}
\usepackage[utf8]{inputenc}
\usepackage{lmodern}
\usepackage[margin=1in]{geometry}
\usepackage{amsmath,amssymb,amsthm,mathtools}
\usepackage[colorlinks=true,linkcolor=blue,citecolor=blue,urlcolor=blue]{hyperref}
\allowdisplaybreaks
\newtheorem{theorem}{Theorem}
\newtheorem{proposition}[theorem]{Proposition}
\newtheorem{lemma}[theorem]{Lemma}
\newtheorem{corollary}[theorem]{Corollary}
\theoremstyle{definition}
\newtheorem{definition}[theorem]{Definition}
\newcommand{\R}{\mathbb R}
\newcommand{\E}{\mathbb E}
\newcommand{\Prb}{\mathbb P}
\newcommand{\Q}{\mathbb Q}
\newcommand{\Ind}{\mathbf 1}
\DeclareMathOperator{\SC}{SC}
\DeclareMathOperator{\OPT}{OPT}
\DeclareMathOperator{\conv}{conv}
\DeclareMathOperator{\supp}{supp}
\title{Envy-free lotteries: a separating certificate and a proportional-component construction}
\author{GPT 6 Astra Ultra}
\date{9 October 2026}
\begin{document}
\maketitle
\noindent\textbf{Problem ID:} \texttt{D63-158}. \textbf{Status: unresolved.}

\section{Target and scope}
For at least three agents with nonnegative additive values, the target asks
for a lottery that is envy-free in expectation and whose every supported
allocation is EF1 and deterministically Pareto optimal. Aziz
\cite[Section 7]{Aziz} leaves the general case open. We give an exact finite
alternative identifying what a counterexample must certify, and a constructive
positive result for independent components with proportional values.
The construction actually has all-good EFX in every outcome. No claim
of novelty is made for the proportional special case or finite separation.

Let $\mathcal G$ be the finite set of allocations that are both EF1 and
Pareto optimal. For $i\ne j$ define
\[
 d_{ij}(A)=v_i(A_i)-v_i(A_j),\qquad d(A)\in\R^{q},\quad q=n(n-1).
\]
Checking membership in $\mathcal G$ by enumerating all complete allocations
is finite, although generally exponential. The following result is a
certificate reduction, not a polynomial-time algorithm.

\section{An exact alternative and a small-support theorem}
\begin{theorem}[Finite separation]\label{alternative}
Suppose $\mathcal G\ne\varnothing$. Exactly one of the following holds:
\begin{enumerate}
\item there is a lottery on $\mathcal G$ with $\E d_{ij}(A)\ge0$ for all $i\ne j$;
\item there are weights $w_{ij}\ge0$, not all zero, and $\gamma>0$ such that
\[
 \sum_{i\ne j}w_{ij}d_{ij}(A)\le-\gamma
 \quad\text{for every }A\in\mathcal G.
\]
\end{enumerate}
Whenever the first alternative holds it has a solution supported on at
most $q+1$ allocations. For rational valuations a rational such lottery
exists. If $\mathcal G=\varnothing$, no lottery exists trivially.
\end{theorem}
\begin{proof}
Let $K=\conv\{d(A):A\in\mathcal G\}$ and let $O=\R^q_{\ge0}$.
A desired lottery exists precisely when $K\cap O\ne\varnothing$.
A vector with the inequalities in the second alternative has nonnegative
dot product with every point in $O$ but negative dot product with every
point in $K$, so the alternatives cannot both hold.

If they are disjoint, minimize $\sum_\ell(\min\{x_\ell,0\})^2$ over
compact $K$, and let $x^*$ be a minimizer. Set $y^*=(x^*)^+$ and
$w=y^*-x^*\ge0$. Disjointness gives $w\ne0$. The objective is continuously
differentiable with gradient $2(x-y)$ at $x=x^*$, where $y=x^+$.
For every $z\in K$, the directional derivative along the segment toward
$z$ is nonnegative, so
$-w\cdot(z-x^*)\ge0$. As $w_\ell>0$ only where $x^*_\ell<0$,
$w\cdot x^*=-\|w\|_2^2$. Therefore
$w\cdot z\le-\|w\|_2^2$, proving the second alternative with
$\gamma=\|w\|_2^2$.

For support size, start with any feasible lottery. If more than $q+1$
weights are positive, the vectors $(1,d(A))$ on its support are linearly
dependent. Choose nonzero coefficients $c_A$ with $\sum c_A=0$ and
$\sum c_A d(A)=0$. The coefficients have both signs. Subtract
$t c_A$ from the weights, where $t=\min_{c_A>0}\lambda_A/c_A$.
Weights stay nonnegative, at least one vanishes, and the total weight and
expected envy vector stay unchanged. Iterate.

For rational input, the original feasible set of lottery weights is a
nonempty bounded polytope defined by rational inequalities. It has a
vertex: minimize the coordinates successively over the polytope, retaining
the minimizer face each time; after all coordinates the remaining face is
a point. At a vertex a full-rank set of active constraints determines the
point, which is rational by Gaussian elimination. Start the preceding
support elimination with this rational lottery and use a rational linear
dependence at every step. All resulting weights remain rational.
\end{proof}
The useful feature is the quantifier over \emph{all} EF1--PO outcomes.
Failure of a proposed rounding method, or a list of individually envious
outcomes, does not constitute the second alternative.

\section{A constructive positive class}
Partition the agents into nonempty components $N_1,\ldots,N_r$ and the
goods into $M_0,M_1,\ldots,M_r$. Every good in $M_0$ has value zero to
everyone. For every $\ell\ge1$ assume weights $b_g\ge0$ on $M_\ell$ and
positive scales $c_i$ for $i\in N_\ell$ such that
\[
 a_{ig}=c_i b_g\quad(i\in N_\ell,g\in M_\ell),\qquad
 a_{ig}=0\quad(i\notin N_\ell,g\in M_\ell).
\]
Zero weights within $M_\ell$ may be moved to $M_0$. Thus every remaining
good in a nonzero component is valued positively by all its agents.

Within component $\ell$, sort its goods by nonincreasing $b_g$ and place
each in a currently least-loaded one of $|N_\ell|$ initially empty bundles,
where load means the sum of $b_g$. Break ties deterministically.
Independently in each component, uniformly permute these bundles among
its agents. Finally put all globally zero goods in any one bundle of
minimum load in its component, or, equivalently, include them in that
component's greedy construction before permuting. If there are no goods
outside $M_0$, choose any component and perform this last operation there.

\begin{lemma}[Descending greedy partitions]\label{greedy}
For any nonnegative item weights, the descending least-load rule produces
bundles $B_1,\ldots,B_k$ satisfying
\[
 b(B_i)\ge b(B_j\setminus\{g\})\quad
 \text{for all }i,j\text{ and every }g\in B_j.
\]
\end{lemma}
\begin{proof}
If $B_j$ is empty there is nothing to check. Let $g_*$ be its last assigned
good. Just before that assignment its load was at most the current load
of every other bundle, hence at most every other bundle's final load.
Because the goods were processed in nonincreasing order, $b_{g_*}\le b_g$
for every $g\in B_j$. Thus
$b(B_j\setminus\{g\})\le b(B_j)-b_{g_*}\le b(B_i)$.
This also handles zero goods: if its last good has weight zero, that bundle's
final load is itself at most every other final load.
\end{proof}
\begin{theorem}[Proportional-component lottery]\label{positive}
Every instance of the displayed component form admits a lottery that is
ex ante envy-free and ex post Pareto optimal and all-good EFX, hence EF1.
For rational input, its ex ante marginal assignment
probabilities can be computed in polynomial time, and an outcome can be
sampled in expected polynomial time using unbiased random bits. It also has a lottery
representation with at most $n(n-1)+1$ supported outcomes.
\end{theorem}
\begin{proof}
Fix an agent $i\in N_\ell$. Against another agent in $N_\ell$,
Lemma~\ref{greedy}, multiplied by $c_i$, gives all-good EFX for every
permutation. Goods from other components are valued zero by $i$, so her
comparison with their bundles has right side zero. Globally zero goods
were added through the greedy zero-weight step in one component; this
preserves the lemma's all-good conclusion there and affects no other
valuation. Thus every outcome is EFX and EF1.

Under a uniform component permutation, each component agent receives each
of its bundles with probability $1/|N_\ell|$. Agent $i$ values her own
expected bundle and every other member's expected bundle at the same
number $c_i b(M_\ell)/|N_\ell|$. Bundles outside her component have
expected value zero. Hence all expected envy inequalities hold.

For efficiency assign agent $i$ welfare weight $1/c_i>0$. Each good in
$M_\ell$ contributes at most $b_g$ to the weighted welfare of any complete
allocation, since outside-component holders value it zero; goods in $M_0$
contribute zero. Every constructed outcome attains the sum of these upper
bounds because goods are assigned within their component. A Pareto
improvement would strictly increase this positively weighted welfare,
which is impossible. This proves deterministic Pareto optimality.

Sorting and rational load comparisons have polynomial cost. A uniform
random permutation can be sampled in expected polynomial bit time by
successive uniform choices with rejection sampling. Each good in a component is marginally assigned uniformly
within that component, giving the marginal probabilities directly.
Theorem~\ref{alternative} gives the small-support existence assertion;
we do not claim that enumerating all support outcomes for its compression
is polynomial in the input size.
\end{proof}
The efficiency proof uses a common supporting welfare functional.
For general heterogeneous valuations such a functional need not be
compatible with every fairness-preserving permutation.

\section{The unresolved step}
The full problem would follow if, for arbitrary additive valuations,
every nonnegative array $w$ admitted an EF1--PO outcome with
$\sum w_{ij}d_{ij}(A)\ge0$. Theorem~\ref{alternative} makes that statement
exactly sufficient and necessary when the good-outcome set is nonempty.
The proportional-component construction proves it only on the stated
restricted domain. No general inequality establishing it, and no negative
certificate contradicting it, is obtained here.

The cited HTML source was checked on 9 October 2026. Its Section 7 expressly
distinguishes the unresolved $n\ge3$ EF1 question from stronger EFX or
fractional-efficiency variants. The proof here uses deterministic PO as
the target does. No finite experiment is used to infer existence.
\begin{thebibliography}{9}
\bibitem{Aziz} H. Aziz, \emph{Best-of-Both-Worlds Fairness and Pareto
Optimality}, arXiv:2608.08966v1, 2026. Section 7.
\url{https://arxiv.org/html/2608.08966v1}.
\end{thebibliography}

\end{document}
